\lecture{3}{Tue 11 Feb 2025 12:30}{Cauchy-Riemann}
There are a few Cauchy-Riemann conditions. I think we are viewing a complex function $f(x+iy)$ as a real function $f(x,y)$.
\begin{itemize}
	\item
		\[
			i \frac{\partial f}{\partial x} =\frac{\partial f}{\partial y} 
		.\]
	\item $f(x,y)=u(x,y)+iv(x,y)$.
	\item 
		\[
			-\frac{\partial v}{\partial x} =\frac{\partial u}{\partial x} 
		.\]
	\item 
		\[
			\frac{\partial u}{\partial x} =\frac{\partial v}{\partial y} 
		.\]
\end{itemize}
If this is true, this means $f$ is a function of $z$ only and not $z^*$. If these conditions are satisfied, $f'(z)$ exists.
\begin{definition}[Analytic]\label{dfn:1.1}
	A function $f(z)$ is analytic in an open region $R\subseteq \mathbb{C} $ if and only if $f'(z)$ exists in $R$. $f$ is entire if $R=\mathbb{C} $.
\end{definition}

For example, $f(z)=z^2$ is entire because $f'(z)=2z$. Similarly, $f(z)=e^z$ is entire because $f'(z)=f(z)$. Given an analytic function, we can talk about singularities.
\begin{definition}[Singularity]\label{dfn:1.2}
	A singularity of an analytic function is an isolated point $z$ such that $f$ is not analytic at $z$, but is analytic in an open disc surrounding $z$.
\end{definition}

For example, $f(z)=1/z$ has a singularity at $z=0$. We often call this a \emph{pole}. Similarly, $f(z)=1/z^2$ has a singularity at $z=0$ as well. More generally,
\[
	f(z)=\frac{R(z_0)}{z-z_0} 
\]
has a singularity at $z_0$. We call $R(z_0)$ the \emph{residue} of the pole, and $z_0$ is the location of the pole.

For example, the residue of $e^z/z$ is
\[
	\lim_{z \to 0} zf(z)=1
.\]

For example,
\[
	f(z)=\frac{1}{z+i} =\frac{1}{z-i} =\frac{2z}{z^2+1} 
\]
has poles at $\pm i$.

Poles generalize to higher order poles. Let
\[
	f(z)=\frac{C}{\left( z-z_0 \right) ^n} 
.\]
This function has a pole at $z_0$ of order $n$, with residue $C$.

Another type of singularity is an essential singularity. Consider $f(z)=e^{1/z} $. The taylor series of this is
\[
	f(z)=\sum_{n=0}^{\infty} \frac{1}{n!} \frac{1}{z^n} 
.\]
Note that
\[
	\lim_{z \to 0} z^p f(z)
\]
is not finite for all $p\in\mathbb{N} $.

The final type of singularity we will discuss is the branch cut. For example, consider $f(z)=\sqrt{z} $. There is some ambiguity in the square root, so branch cuts will sort this out. If we write
\[
	z=\exp i\theta 
,\]
then
\[
	z^{1/2} =\exp (i\theta /2)
.\]
However, we could have also written
\[
	z=\exp (i\left( \theta +2\pi  \right) )
.\]
We require this phase angle to be the same, otherwise we have a branch cut.

A contour integral is an integral from $z_1$ to $z_2$, where we integrate along a path $P$ from $z_1$ to $z_2$. That is, we take the integral
\[
	\int_{P} f(z) \,\mathrm{d} z
.\]
Computing this is the same as a multivariate limit; just parameterize $P$ with a function $z(t)=x(t)+iy(t)$ for $t\in[0,1]$, and then take the integral
\[
	\int_{0}^{1} f(z(t)) \frac{\mathrm{d}z}{\mathrm{d}t} \,\mathrm{d} t
.\]

Path independence is when two distinct paths $P_1,P_2$ have the same integral
\[
	\int_{P_1} f(z) \,\mathrm{d} z=\int_{P_2} f(z) \,\mathrm{d} z
.\]

